\textbf{Proof.}
Fix a pilot realization and an endpoint \(e\in\{L,U\}\).  The empirical matrix \(\widehat Q_e\) is positive semidefinite: draw the \(n\) signed potential-outcome pairs independently from the empirical comonotone coupling when \(e=U\), or from the empirical countermonotone coupling when \(e=L\).  The block calculation in def:moment-matrix shows that \(\widehat Q_e\) is the second-moment matrix of the resulting signed \(2n\)-vector.  Hence
\[
\widehat A_e
=\Omega^{1/2}\Pi^{-1}\widehat Q_e\Pi^{-1}\Omega^{1/2}
\succeq0.
\tag{1}
\]
Assumption ass:exposure-positivity supplies the inverse of \(\Pi\).

For a rank-two class, \(S_B=B^{\top}\Omega B\succ0\), inversion is continuous, and the displayed formula in def:pilot-plugin shows that \(B\mapsto\widehat w_m^+(B)\) is continuous on every normalized rank-two quotient chart.  We next adapt the one-sided density argument of lem:rank-stratified-projector-reduction to the empirical endpoint matrices.  Let \([B]_+\) have rank at most one.  Since \(\mathcal B_Z\ne\varnothing\), choose a rank-two representative \(C^\circ\).  The perturbation and normalization construction in that lemma gives normalized rank-two classes \([B_k]_+\) and a subsequential projector limit \(P_\infty\) such that
\[
P_\infty-P_B\succeq0.
\tag{2}
\]
The projector trace identity, with \(Q(c_e)\) replaced by \(\widehat Q_e\), is
\[
n^2\widehat g_{B,e}^+
=\operatorname{tr}(P_B\widehat A_e).
\tag{3}
\]
Equations (1)--(3) imply
\[
\lim_k\widehat g_{B_k,e}^+
=n^{-2}\operatorname{tr}(P_\infty\widehat A_e)
\ge n^{-2}\operatorname{tr}(P_B\widehat A_e)
=\widehat g_{B,e}^+.
\tag{4}
\]
The uncontrolled empirical endpoint term is independent of the basis.  Subtracting (4) from that term and taking the maximum over the two endpoints yields
\[
\limsup_{k\to\infty}\widehat w_m^+(B_k)
\le\widehat w_m^+(B).
\tag{5}
\]

The rational directions used in def:pilot-plugin are dense on each normalized sphere.  Positive definiteness of \(C^{\top}\Omega C\) is open, so the retained sequence \(\mathcal D_Z\) is dense in every rank-two quotient chart.  By rank-two continuity, for each \(k\) there is an enumerated class \([C_{j_k}]_+\) satisfying
\[
\widehat w_m^+(C_{j_k})
\le\widehat w_m^+(B_k)+k^{-1}.
\tag{6}
\]
Equations (5)--(6) show, for every singular class,
\[
\inf_{j\ge1}\widehat w_m^+(C_j)
\le\widehat w_m^+(B).
\tag{7}
\]
The same inequality for a rank-two class follows directly from density and continuity.  Taking the infimum of the right side over the full carrier gives
\[
\inf_{j\ge1}\widehat w_m^+(C_j)
\le\inf_{[B]_+\in\mathcal B_Z^+}\widehat w_m^+(B).
\tag{8}
\]
The reverse inequality holds because every \([C_j]_+\) belongs to \(\mathcal B_Z^+\).  This proves the infimum equality without assuming attainment.

It remains to prove measurability.  For fixed \(j\), the empirical means, variances, and endpoint quantile products are Borel functions of the finitely many pilot observations; equivalently, they are finite sums and order-statistic functions.  Since \(S_{C_j}\succ0\) is fixed by the design, the ordinary inverse in \(\widehat w_m^+(C_j)\) is deterministic.  Thus
\[
X_j:=\widehat w_m^+(C_j)
\]
is pilot measurable.  The countable infimum \(I_m=\inf_{j\ge1}X_j\) is measurable.  Strict positivity of \(\varepsilon_m\) and the defining property of the infimum imply that the qualifying set \(\{j:X_j\le I_m+\varepsilon_m\}\) is nonempty.  Therefore its least index is finite, and
\[
\{J_{m,\varepsilon_m}=j\}
=\{X_j\le I_m+\varepsilon_m\}
\cap\bigcap_{k<j}\{X_k>I_m+\varepsilon_m\}
\tag{9}
\]
is measurable.  A countable-valued map with the measurable fibers (9) is measurable, so \([C_{J_{m,\varepsilon_m}}]_+\) is a measurable quotient-valued selector.  Its defining inequality and the proved infimum equality give the final display. \(\square\)